
\begin{coro}\label{cor: main thm of section 2}
For any nondegenerate quiver with potential $(Q,W)$ of finite type, there is a $1$-$1$ correspondence between maximal green sequences for $Q$ and isomorphism classes of complete FHO sequences $M_1,\ldots,M_m$ in $J(Q,W)\text{-}\modr $ where the green $c$-sequence corresponding to $(M_i)$ is $(\beta_i=\undim M_i)$. \qed
\end{coro}

\begin{coro}[Rotation Lemma]\label{cor: rotation lemma}
There is a $1$-$1$ correspondence between complete FHO sequences $M_1,\ldots,M_m$ in $J(Q,W)$ starting with $M_1=S_k$ and complete FHO sequences $X_1,\ldots,X_m$ of the same length for $J(\mu_k(Q,W))$ ending with $X_m=S_k'$. The correspondence is given by $X_i=\psi_k(M_{i+1})$ for $i<m$.
\end{coro}

\begin{proof} The correspondence is already given by $X_i=\psi_k(M_{i+1})$ for $i<m$ and $M_j=\psi_k^{-1}(X_{j-1})$ for $j>1$. It remains to show that this formula sends complete FHO sequences in $\Lambda\text{-}\modr $ to complete FHO sequences in $\Lambda'\text{-}\modr $ and vice versa. But this follows from Proposition~\ref{prop: complete iff M1 is simple and M2, etc is maximal} and the fact that $\psi_k:S_k^\perp\cong \,^\perp S_k'$ is an equivalence of categories: $M_1,\ldots,M_m$, with $M_1=S_k$ and $M_2,\ldots,M_m\in M_1^\perp$ is a complete FHO sequence for $\Lambda\text{-}\modr $ if and only if $M_2,\ldots,M_m$ is a maximal weakly FHO sequence in $S_k^\perp$. Since $\psi_k$ is an equivalence, this is equivalent to $X_1,\ldots,X_{m-1}$ being a maximal weakly FHO sequence in $\,^\perp S_k'$ in $\Lambda'\text{-}\modr $. By~\ref{prop: complete iff M1 is simple and M2, etc is maximal} this is equivalent to $X_1,\ldots,X_{m-1},S_k'$ being a complete FHO sequence in $\Lambda'\text{-}\modr $.
\end{proof}
